\section{Introduction}
\label{sec:intro}

Flow-based intrusion detectors routinely report accuracies above 99\% on benchmarks such as
CIC-IDS2017~\cite{sharafaldin2018}, mostly with random train/test splits in which flows of one attack session
fall on both sides. Network traffic is time-dependent, and estimates from a chronological split differ from those
of a random split~\cite{nismara2026temporal}; careful sequence evaluations find that the protocol matters far more
than the model architecture~\cite{moczkodan2026transformers}.

A deployed detector faces three further conditions. \emph{New attack classes} appear after training, and a
closed-set model assigns them to a known class, most often benign. \emph{Labels are scarce}: an analyst inspects
a small fraction of the flows. \emph{Labels are late}: ground truth arrives after the decisions were made. In
addition, every model update is a deployment, and a bad one degrades detection silently. Prior work handles
these conditions one at a time (\cref{sec:related}).

We study them together and propose \nova{} (Novelty-aware, Open-world, Validated Adaptation). Our
contributions are:
\begin{enumerate}
  \item a protocol for drift-aware open-world NIDS: chronological replay in windows, a label oracle that enforces
        one budget and delay for every method, and metrics that stay comparable when methods learn different
        classes (\cref{sec:problem});
  \item \nova{}, which combines class-conditional $k$-NN rejection with per-class recalibration, a
        novelty-weighted, cluster-diverse query, agreement-checked label propagation, a class-balanced memory
        and a guarded deployment with rollback (\cref{sec:method});
  \item a comparison with eight methods, including four re-implemented from the literature, on three datasets,
        with an ablation of every component (\cref{sec:setup,sec:results}). Code, configurations and per-flow
        predictions will be released with the final version.
\end{enumerate}
